\documentclass{ximera}

\input{../../../preamble.tex}


\definecolor{fvdnavy}{HTML}{071D43}
\definecolor{fvdcyan}{HTML}{20BFC6}
\definecolor{fvdcyanlight}{HTML}{EFFCFB}
\definecolor{fvdmagenta}{HTML}{F10D69}
\definecolor{fvdmagentalight}{HTML}{FFF1F7}
\definecolor{fvdtext}{HTML}{163044}





%=================================================
% Pedagogische omgevingen
% PDF: tcolorbox
% HTML: learning-box
%=================================================

\newenvironment{voorkennis}
{%
    \ifdefined\HCode
        \HCode{%
<section class="learning-box learning-box--prior">
<div class="learning-box__header">
<span class="learning-box__icon" aria-hidden="true"></span>
<span class="learning-box__title">Voorkennis</span>
</div>
<div class="learning-box__content">
        }%
        \begin{itemize}
    \else
       \begin{tcolorbox}[
    enhanced,
    breakable,
    colback=fvdcyanlight,
    colframe=fvdcyan,
    coltext=fvdtext,
    boxrule=0.5pt,
    leftrule=4pt,
    arc=4mm,
    outer arc=4mm,
    left=7mm,
    right=6mm,
    top=5mm,
    bottom=5mm,
    before skip=6mm,
    after skip=6mm,
    title=Voorkennis,
    coltitle=fvdcyan,
    fonttitle=\bfseries\Large,
    attach title to upper={\par\smallskip},
    boxed title style={
        empty,
        left=0mm,
        right=0mm,
        top=0mm,
        bottom=0mm
    },
    overlay={
        \draw[
            fvdcyan,
            line width=0.5pt
        ]
        ([xshift=42mm,yshift=-13mm]frame.north west)
        --
        ([xshift=-7mm,yshift=-13mm]frame.north east);
    }
]
        \begin{itemize}
    \fi
}
{%
    \end{itemize}
    \ifdefined\HCode
        \HCode{%
</div>
</section>
        }%
    \else
        \end{tcolorbox}
    \fi
}


\newenvironment{leerdoelen}
{%
    \ifdefined\HCode
        \HCode{%
<section class="learning-box learning-box--goals">
<div class="learning-box__header">
<span class="learning-box__icon" aria-hidden="true"></span>
<span class="learning-box__title">Leerdoelen</span>
</div>
<div class="learning-box__content">
        }%
        \begin{itemize}
    \else
\begin{tcolorbox}[
    enhanced,
    breakable,
    colback=fvdmagentalight,
    colframe=fvdmagenta,
    coltext=fvdtext,
    boxrule=0.5pt,
    leftrule=4pt,
    arc=4mm,
    outer arc=4mm,
    left=7mm,
    right=6mm,
    top=5mm,
    bottom=5mm,
    before skip=6mm,
    after skip=6mm,
    title=Leerdoelen,
    coltitle=fvdmagenta,
    fonttitle=\bfseries\Large,
    attach title to upper={\par\smallskip},
    boxed title style={
        empty,
        left=0mm,
        right=0mm,
        top=0mm,
        bottom=0mm
    },
    overlay={
        \draw[
            fvdmagenta,
            line width=0.5pt
        ]
        ([xshift=42mm,yshift=-13mm]frame.north west)
        --
        ([xshift=-7mm,yshift=-13mm]frame.north east);
    }
]
        \begin{itemize}
    \fi
}
{%
    \end{itemize}
    \ifdefined\HCode
        \HCode{%
</div>
</section>
        }%
    \else
        \end{tcolorbox}
    \fi
}


\addPrintStyle{../../..}

\title{Correlatie --- Hoe sterk is het lineaire verband?}
\author{Ingmar Herreman}

\begin{abstract}
In het vorige hoofdstuk leerden we een puntenwolk systematisch lezen:
we onderzochten de richting, de vorm, de sterkte en bijzondere
kenmerken van het patroon.

Wanneer een verband ongeveer lineair is, willen we die visuele
beschrijving nauwkeuriger maken. Daarvoor gebruiken we de
\textbf{correlatiecoëfficiënt} $r$.

De correlatiecoëfficiënt vat twee eigenschappen van een lineair
verband samen: de \textbf{richting} en de \textbf{sterkte}.
Maar één getal vertelt nooit het volledige verhaal. Daarom leren we
$r$ steeds samen met het spreidingsdiagram en de context interpreteren.

We onderzoeken ook situaties waarin $r$ misleidend kan zijn:
niet-lineaire verbanden, uitschieters en datasets met een bijzondere
structuur. Zo leren we niet alleen een correlatiecoëfficiënt aflezen,
maar vooral beoordelen wat die coëfficiënt wel en niet over de
gegevens vertelt.
\end{abstract}

\begin{document}

% FVD_CHAPTER_DATA_START
% Automatisch gegenereerd uit index.tex.
% Niet handmatig aanpassen.
\fvdchapterdata{3}{Verbanden onderzoeken met data}{6}
% FVD_CHAPTER_DATA_END



\maketitle

\begin{voorkennis}
  \item Je weet wat gekoppelde numerieke gegevens zijn.
  \item Je kan een spreidingsdiagram maken en lezen.
  \item Je kan de richting, vorm en visuele sterkte van een puntenwolk beschrijven.
  \item Je kan een ongeveer lineair verband onderscheiden van een duidelijk niet-lineair verband.
  \item Je kan uitschieters en clusters herkennen.
\end{voorkennis}

\begin{leerdoelen}
  \item Je kan uitleggen waarom een visuele beoordeling van een lineair verband niet altijd volstaat.
  \item Je kent het bereik $-1\leq r\leq 1$ van de correlatiecoëfficiënt.
  \item Je kan het teken van $r$ koppelen aan de richting van een lineair verband.
  \item Je kan $|r|$ koppelen aan de sterkte van een lineair verband.
  \item Je kan een correlatiecoëfficiënt steeds samen met een spreidingsdiagram interpreteren.
  \item Je kan uitleggen waarom $r\approx 0$ niet noodzakelijk betekent dat er geen verband bestaat.
  \item Je kan onderzoeken hoe een uitschieter de correlatiecoëfficiënt kan beïnvloeden.
  \item Je kan met ICT de correlatiecoëfficiënt van een dataset bepalen.
  \item Je kan op basis van puntenwolk, correlatiecoëfficiënt en context een onderbouwde conclusie formuleren.
\end{leerdoelen}


% ============================================================
\FVDsection{Van kijken naar meten}
% ============================================================

In het vorige hoofdstuk beschreven we een puntenwolk met woorden.

Bijvoorbeeld:

\begin{quote}
De punten vertonen een sterk positief en ongeveer lineair verband.
\end{quote}

Zo'n beschrijving is nuttig, maar bevat ook een subjectief element.
Wat de ene leerling ``sterk'' noemt, kan een andere leerling
``matig sterk'' noemen.

Wanneer een puntenwolk ongeveer lineair is, willen we daarom een
numerieke maat voor de richting en de sterkte van dat lineaire verband.

\begin{opmerking}
De volgorde blijft belangrijk:

\[
\important{
\text{puntenwolk bekijken}
\longrightarrow
\text{vorm beoordelen}
\longrightarrow
\text{correlatie interpreteren}
}
\]

We berekenen dus niet blind een getal. Eerst onderzoeken we of een
lineaire beschrijving van de gegevens zinvol lijkt.
\end{opmerking}


% ============================================================
\FVDsection{De correlatiecoëfficiënt}
% ============================================================

\begin{definitie}
De \textbf{correlatiecoëfficiënt} $r$ is een getal dat de
\textbf{richting} en de \textbf{sterkte van het lineaire verband}
tussen twee numerieke grootheden beschrijft.

Voor $r$ geldt steeds

\[
\boxed{-1\leq r\leq 1}.
\]
\end{definitie}

De correlatiecoëfficiënt combineert dus twee soorten informatie:

\[
\important{
\text{teken van }r
\longrightarrow
\text{richting}
}
\]

en

\[
\important{
|r|
\longrightarrow
\text{sterkte van het lineaire verband}
}
\]


% ------------------------------------------------------------
\FVDsubsection{Het teken van r}
% ------------------------------------------------------------

Wanneer

\[
r>0,
\]

is de lineaire samenhang \textbf{positief}: grotere waarden van de ene
grootheid gaan doorgaans samen met grotere waarden van de andere.

Wanneer

\[
r<0,
\]

is de lineaire samenhang \textbf{negatief}: grotere waarden van de ene
grootheid gaan doorgaans samen met kleinere waarden van de andere.

Het teken zegt niets over de sterkte.

Zo kunnen

\[
r=0{,}20
\qquad\text{en}\qquad
r=0{,}92
\]

beide positief zijn, terwijl het tweede lineaire verband veel sterker is.


% ------------------------------------------------------------
\FVDsubsection{De absolute waarde van r}
% ------------------------------------------------------------

Voor de sterkte kijken we naar

\[
|r|.
\]

Hoe dichter $|r|$ bij $1$ ligt, hoe dichter de punten bij een rechte
band liggen.

Hoe dichter $|r|$ bij $0$ ligt, hoe zwakker de \textbf{lineaire}
samenhang.

\[
\begin{array}{c|c}
|r| & \text{globale interpretatie}\\ \hline
\text{dicht bij }1 & \text{sterk lineair verband}\\
\text{dicht bij }0 & \text{zwak lineair verband}
\end{array}
\]

\begin{opmerking}
Er bestaan conventies die woorden zoals ``zwak'', ``matig'' en
``sterk'' aan vaste intervallen voor $|r|$ koppelen.

We behandelen zulke grenzen niet als natuurwetten. De betekenis van
een correlatie hangt ook af van de context, het aantal waarnemingen,
de kwaliteit van de gegevens en de vorm van de puntenwolk.
\end{opmerking}


% ============================================================
\FVDsection{De bijzondere waarden}
% ============================================================

\FVDsubsection{Perfect positieve lineaire samenhang}

Als alle punten exact op een stijgende rechte liggen, dan

\[
r=1.
\]

\begin{opmerking}
$r=1$ betekent niet dat de richtingscoëfficiënt van die rechte gelijk
is aan $1$.

Het betekent dat de punten exact op \emph{een} stijgende rechte liggen.
\end{opmerking}


\FVDsubsection{Perfect negatieve lineaire samenhang}

Als alle punten exact op een dalende rechte liggen, dan

\[
r=-1.
\]

Ook hier zegt $-1$ niets over de numerieke helling van de rechte.
Het geeft aan dat de lineaire samenhang perfect negatief is.


\FVDsubsection{Wat betekent r=0?}

Wanneer

\[
r=0,
\]

is er geen lineaire samenhang in de gegevens zoals die door $r$ wordt
gemeten.

Maar daaruit volgt niet:

\[
\text{``er bestaat geen enkel verband.''}
\]

Een duidelijke U-vormige puntenwolk kan bijvoorbeeld een sterk
niet-lineair verband vertonen terwijl $r$ ongeveer nul is.

\begin{opmerking}
Dit is één van de belangrijkste interpretatieregels van dit hoofdstuk:

\[
\important{
r\approx 0
\quad\not\Rightarrow\quad
\text{geen verband}
}
\]

Wel:

\[
\important{
r\approx 0
\quad\Rightarrow\quad
\text{weinig of geen lineaire samenhang volgens }r
}
\]
\end{opmerking}


% ============================================================
\FVDsection{Puntenwolk en r horen samen}
% ============================================================

De correlatiecoëfficiënt is een samenvatting. Bij een samenvatting gaat
altijd informatie verloren.

Daarom bekijken we nooit alleen $r$.

\begin{voorbeeld}[Vier mogelijke conclusies]
Stel dat vier datasets de volgende correlatiecoëfficiënten hebben:

\[
r_A=0{,}94,\qquad
r_B=-0{,}91,\qquad
r_C=0{,}18,\qquad
r_D=-0{,}35.
\]

Dan kunnen we voorlopig zeggen:

\begin{itemize}
  \item dataset $A$: sterke positieve lineaire samenhang;
  \item dataset $B$: sterke negatieve lineaire samenhang;
  \item dataset $C$: zwakke positieve lineaire samenhang;
  \item dataset $D$: eerder zwakke negatieve lineaire samenhang.
\end{itemize}

Maar vóór een definitieve conclusie bekijken we ook de puntenwolken.
Misschien is er een uitschieter, een krom patroon of een opsplitsing in
clusters.
\end{voorbeeld}


\begin{oefening}[Lees r correct]{\oefB \ESS}
\begin{question}
Wat vertelt $r=-0{,}88$?

\begin{multipleChoice}
  \choice{Een sterk positief lineair verband.}
  \choice[correct]{Een sterk negatief lineair verband.}
  \choice{Een zwak negatief lineair verband.}
  \choice{Een sterk niet-lineair verband.}
\end{multipleChoice}
\end{question}

\begin{question}
Welke correlatie is lineair het sterkst?

\[
r_1=-0{,}93,\qquad r_2=0{,}81,\qquad r_3=0{,}42
\]

\begin{multipleChoice}
  \choice[correct]{$r_1$}
  \choice{$r_2$}
  \choice{$r_3$}
\end{multipleChoice}
\end{question}

\begin{question}
Waarom vergelijk je daarvoor de absolute waarden?

\begin{oplossing}
Omdat het teken de richting aangeeft. De sterkte van de lineaire
samenhang wordt beschreven door de grootte van $|r|$.
\end{oplossing}
\end{question}
\end{oefening}


% ============================================================
\FVDsection{Wat r niet ziet}
% ============================================================

Een correlatiecoëfficiënt kan correct berekend zijn en toch een
onvolledig beeld geven.

Daarom onderzoeken we drie belangrijke valkuilen.


% ------------------------------------------------------------
\FVDsubsection{Niet-lineaire verbanden}
% ------------------------------------------------------------

Beschouw opnieuw een U-vormig verband zoals

\[
y\approx x^2.
\]

Voor negatieve $x$ daalt $y$ wanneer $x$ toeneemt.
Voor positieve $x$ stijgt $y$ wanneer $x$ toeneemt.

Over de volledige dataset kunnen die twee richtingen elkaar in de
lineaire correlatie grotendeels opheffen.

Toch is het patroon duidelijk.

\begin{opmerking}
Een lage waarde van $|r|$ kan dus twee heel verschillende situaties
verbergen:

\begin{enumerate}
  \item er is werkelijk nauwelijks een patroon;
  \item er is een duidelijk patroon, maar het is niet lineair.
\end{enumerate}

Alleen de puntenwolk laat ons dat onderscheid zien.
\end{opmerking}


% ------------------------------------------------------------
\FVDsubsection{Uitschieters}
% ------------------------------------------------------------

Een uitzonderlijk punt kan $r$ sterk beïnvloeden.

Soms maakt één punt een lineair verband schijnbaar sterker.
In andere datasets maakt één punt een bestaand lineair verband juist
veel zwakker.

Daarom vragen we bij een uitschieter:

\begin{itemize}
  \item is het punt correct gemeten en ingevoerd?
  \item hoort het bij dezelfde populatie of experimentele situatie?
  \item hoe verandert de puntenwolk als we het punt onderzoeken?
  \item hoe verandert $r$?
\end{itemize}

\begin{opmerking}
Een uitschieter verwijderen omdat hij ``niet mooi past'' is geen
verantwoorde statistische werkwijze.
\end{opmerking}


% ------------------------------------------------------------
\FVDsubsection{Clusters}
% ------------------------------------------------------------

Ook clusters kunnen een globale correlatie misleidend maken.

Stel dat gegevens van twee verschillende materialen samen in één
dataset staan. Binnen elk materiaal kan nauwelijks een verband bestaan,
terwijl de twee groepen samen toch een duidelijke globale richting
lijken te vormen.

Daarom blijft de structuur van de puntenwolk belangrijk.


% ============================================================
\FVDsection{Correlatie bepalen met ICT}
% ============================================================

Voor realistische datasets gebruiken we ICT om $r$ te bepalen.

Het doel is niet om een lange rekenprocedure uit het hoofd uit te
voeren, maar om de uitvoer correct te interpreteren.

Een goede werkwijze is:

\begin{enumerate}
  \item controleer de gekoppelde gegevens;
  \item maak eerst een spreidingsdiagram;
  \item beoordeel of het verband ongeveer lineair lijkt;
  \item laat ICT de correlatiecoëfficiënt $r$ bepalen;
  \item interpreteer teken en grootte van $r$;
  \item vergelijk die interpretatie met de puntenwolk;
  \item formuleer een conclusie in de context.
\end{enumerate}

\begin{opmerking}
Een antwoord zoals

\[
r=0{,}87
\]

is geen volledige analyse.

Een betere conclusie is bijvoorbeeld:

\begin{quote}
De puntenwolk vertoont een sterk positief en ongeveer lineair verband.
Dat wordt ondersteund door $r=0{,}87$. In deze dataset gaan grotere
waarden van $x$ doorgaans samen met grotere waarden van $y$.
\end{quote}
\end{opmerking}


% ============================================================
\FVDsection{Correlatie is nog geen causaliteit}
% ============================================================

Een grote waarde van $|r|$ zegt dat twee grootheden sterk
\textbf{lineair samenhangen} in de dataset.

Ze zegt niet waarom dat verband bestaat.

Dus zelfs bij

\[
|r|\approx 1
\]

mogen we niet automatisch besluiten dat de ene grootheid de andere
veroorzaakt.

\begin{voorbeeld}[Temperatuur en ijsverkoop]
Stel dat temperatuur en ijsverkoop sterk positief gecorreleerd zijn.

Dat betekent dat warme dagen in de onderzochte gegevens doorgaans
samengaan met een hogere ijsverkoop.

De correlatiecoëfficiënt vertelt ons echter niet op zichzelf welke
causale verklaring correct is. Daarvoor hebben we kennis van de context
en eventueel verder onderzoek nodig.
\end{voorbeeld}

In een later hoofdstuk onderzoeken we dit onderscheid uitgebreid.


% ============================================================
\FVDsection{Een volledige correlatieanalyse}
% ============================================================

Wanneer je een nieuwe dataset krijgt, gebruik je voortaan deze
strategie:

\begin{enumerate}
  \item \textbf{Context:} welke twee numerieke grootheden worden onderzocht?
  \item \textbf{Puntenwolk:} wat zie je qua richting, vorm, sterkte en bijzonderheden?
  \item \textbf{Lineariteit:} is een lineaire beschrijving redelijk?
  \item \textbf{Correlatie:} bepaal $r$ met ICT.
  \item \textbf{Interpretatie:} wat vertellen teken en $|r|$?
  \item \textbf{Controle:} past die numerieke samenvatting bij de puntenwolk?
  \item \textbf{Conclusie:} formuleer een contextgebonden besluit zonder causaliteit te veronderstellen.
\end{enumerate}

\begin{oefening}[Van getal naar conclusie]{\oefU \ESS}
Een onderzoek naar de buitentemperatuur $T$ en het dagelijkse
elektriciteitsverbruik $E$ van een koelinstallatie levert een
ongeveer lineaire stijgende puntenwolk en

\[
r=0{,}91.
\]

\begin{question}
Wat zegt het teken van $r$?

\begin{oplossing}
Het verband is positief: hogere temperaturen gaan in deze dataset
doorgaans samen met een hoger elektriciteitsverbruik.
\end{oplossing}
\end{question}

\begin{question}
Wat zegt $|r|$?

\begin{oplossing}
Omdat $|r|=0{,}91$ dicht bij $1$ ligt, is de lineaire samenhang sterk.
\end{oplossing}
\end{question}

\begin{question}
Formuleer een volledige conclusie.

\begin{oplossing}
De puntenwolk en $r=0{,}91$ wijzen op een sterk positief en ongeveer
lineair verband tussen buitentemperatuur en elektriciteitsverbruik van
de koelinstallatie. Binnen deze dataset gaan hogere temperaturen
doorgaans samen met een hoger verbruik. Uit correlatie alleen volgt
nog geen bewijs van causaliteit.
\end{oplossing}
\end{question}
\end{oefening}


% ============================================================
\FVDsection{Samengevat}
% ============================================================

De correlatiecoëfficiënt $r$ beschrijft de richting en de sterkte van
een \textbf{lineair} verband.

\[
\boxed{-1\leq r\leq 1}
\]

Het teken geeft de richting:

\[
r>0
\Rightarrow
\text{positieve lineaire samenhang},
\]

\[
r<0
\Rightarrow
\text{negatieve lineaire samenhang}.
\]

De absolute waarde geeft informatie over de sterkte:

\[
|r|\text{ dicht bij }1
\Rightarrow
\text{sterke lineaire samenhang},
\]

\[
|r|\text{ dicht bij }0
\Rightarrow
\text{zwakke lineaire samenhang}.
\]

Maar:

\[
\important{
r\approx0
\not\Rightarrow
\text{geen enkel verband}
}
\]

en

\[
\important{
\text{correlatie}
\not\Rightarrow
\text{causaliteit}
}
\]

Daarom gebruiken we altijd:

\[
\boxed{
\text{puntenwolk}
+
r
+
\text{context}
\longrightarrow
\text{conclusie}
}
\]

In het volgende hoofdstuk zetten we een volgende stap:
wanneer het verband voldoende lineair is, proberen we het met een
\textbf{lineair model} te beschrijven.

\end{document}
